problem:  
Let \((M_i, g_i)\) be a sequence of compact \(n\)-dimensional Riemannian manifolds with almost nonnegative sectional curvature,  
\[
\mathrm{sec}_{g_i} \ge -\varepsilon_i, \quad \varepsilon_i \to 0, \quad \mathrm{diam}(M_i) \le 1,
\]
and assume that \((M_i, g_i)\) converge in the Gromov–Hausdorff sense to a compact Alexandrov space \(X\) of curvature \(\ge 0\).  

It is known that locally, collapse under almost nonnegative curvature can often be described by infranil bundles or local submersions with Alexandrov bases, but global limits may exhibit complex singularities exceeding smooth modelings.  
For each such limit space \(X\), define its **minimal smooth embedding dimension**
\[
N_{\min}(X) = \min\{ N \in \mathbb{N} : \exists\ \text{a smooth, complete } (n+N)\text{-manifold } (Y,h) \text{ with } \sec_h \ge 0,
\text{ and a submetry } f\!:Y \to X \}.
\]

**Problem (Minimal smooth submetry dimension for almost nonnegative limits):**  
Is \(N_{\min}(X)\) bounded above by a function of \(n\) only, i.e. does there exist a universal constant \(N(n)\) such that for every Alexandrov space \(X\) arising as a Gromov–Hausdorff limit of \(n\)-dimensional manifolds with almost nonnegative sectional curvature,
\[
N_{\min}(X) \le N(n)?
\]
Furthermore, can one characterize whether \(N_{\min}(X)=0\) (i.e. \(X\) itself admits a smooth nonnegatively curved Riemannian structure) purely in terms of the metric stratification of \(X\)?

Why is it a "good" problem:  
This revised problem formulates a quantitative and intrinsic geometric challenge that goes beyond known results and earlier formulations. It introduces a **minimal-dimensional smooth submetry realization** concept—a natural yet unexplored invariant connecting Alexandrov geometry, Riemannian collapse theory, and transformation group theory.  

Compared with earlier quotient‐realization problems, this one eliminates ambiguities (“metric–quotient‑like”) by invoking the precise notion of a **submetry**, a well‑defined generalization of a Riemannian submersion preserving the metric structure of the limit. The problem also removes the ad hoc fiber restrictions and refocuses on determining the *smallest possible ambient smooth dimension* needed for such a realization.  

If a universal bound \(N(n)\) exists, it would reveal a deep compactness principle for smooth nonnegative curvature in the Gromov–Hausdorff category, showing that all collapse singularities of almost nonnegative limits can be captured within a finite-dimensional smooth curvature world. If not, it would exhibit an essential failure of dimensional uniformity in the smooth resolution of metric limits—a fundamentally new phenomenon.  

The problem bridges several core areas—collapsing theory, Alexandrov geometry, and Riemannian submersions—while remaining simple to state: “Is there a universal smooth ambient size for almost nonnegative curvature limits?” Its clarity, depth, and structural focus make it a genuine research‑level “good” problem.